% Part: first-order-logic
% Chapter: proof-systems
% Section: introduction

\documentclass[../../../include/open-logic-section]{subfiles}

\begin{document}

\iftag{FOL}
      {\olfileid{fol}{prf}{int}}
      {\olfileid{pl}{prf}{int}}

\olsection{પરિચય}

તર્કશાસ્ત્રોમાં સામાન્ય રીતે અર્થવિચાર અને !!a{derivation} તંત્ર બંને
હોય છે. અર્થવિચાર સત્ય, સંતોષ્યતા, પ્રામાણ્ય અને ફલિતતા જેવા ખ્યાલો
સાથે સંબંધ ધરાવે છે. !!{derivation} તંત્રોનો હેતુ ફલિતતા અને
પ્રામાણ્ય સ્થાપવાની સંપૂર્ણપણે વાક્યરચનાત્મક રીત આપવાનો છે. તે એ
અર્થમાં સંપૂર્ણપણે વાક્યરચનાત્મક છે કે આવા તંત્રમાં !!{derivation}
સાન્ત વાક્યરચનાત્મક પદાર્થ હોય છે—સામાન્ય રીતે !!{sentence}s અથવા
!!{formula}sની શ્રેણી (અથવા બીજી કોઈ સાન્ત ગોઠવણી). સારાં
!!{derivation} તંત્રોમાં !!{sentence}s અથવા !!{formula}sની આપેલી કોઈ
પણ શ્રેણી કે ગોઠવણી ``યોગ્ય'' છે કે નહિ તે યાંત્રિક રીતે ચકાસી શકાય છે.

પ્રથમ-ક્રમના તર્કશાસ્ત્ર માટેનાં સૌથી સરળ (અને ઐતિહાસિક રીતે પહેલાં)
!!{derivation} તંત્રો \emph{સ્વયંસિદ્ધિમૂલક} હતાં. આવા તંત્રમાં
!!{formula}sની કોઈ શ્રેણીને !!a{derivation} ગણવા માટે તેમાંનું દરેક
!!{formula} કાં તો ``સ્વયંસિદ્ધિઓ''ના કોઈ નિશ્ચિત ગણમાંનું હોવું જોઈએ,
અથવા શ્રેણીમાં તેની પહેલાં આવતાં !!{formula}sમાંથી નિશ્ચિત સંખ્યાના
``અનુમાન-નિયમો'' પૈકી કોઈ એક વડે મળવું જોઈએ. !!a{formula}
સ્વયંસિદ્ધિ છે કે નહિ અને કોઈ અનુમાન-નિયમથી બીજાં !!{formula}sમાંથી
યોગ્ય રીતે મળે છે કે નહિ તે યાંત્રિક રીતે ચકાસી શકાય છે.
સ્વયંસિદ્ધિમૂલક !!{derivation} તંત્રોનું વર્ણન કરવું અને
અધિસિદ્ધાંતની દૃષ્ટિએ તેમને સંભાળવું સહેલું છે, પરંતુ તેમાંની
!!{derivation}s વાંચવી અને સમજવી અઘરી છે, અને બનાવવી પણ અઘરી છે.

બીજાં !!{derivation} તંત્રો !!{derivation}s રચવાનું અથવા પૂરી થયેલી
!!{derivation}s સમજવાનું સહેલું બને તે હેતુથી વિકસાવાયાં છે. તેના
દાખલા સ્વાભાવિક નિગમન, સત્ય-વૃક્ષો—જેને ટેબ્લો સાબિતીઓ પણ કહે
છે—અને સિક્વન્ટ કલન છે. કેટલાંક !!{derivation} તંત્રો ખાસ કરીને
યાંત્રિક અમલને ધ્યાનમાં રાખીને રચાયાં છે; દાખલા તરીકે, રિઝોલ્યૂશન
પદ્ધતિનો સૉફ્ટવેરમાં અમલ કરવો સહેલો છે (પરંતુ તેની !!{derivation}s
સમજવી લગભગ અશક્ય છે). આવાં મોટા ભાગનાં !!{derivation} તંત્રો
!!{derivation}sને !!{formula}sની
શ્રેણીઓ તરીકે નહિ, પણ તેમનાં વૃક્ષો તરીકે રજૂ કરે છે.
તેથી !!a{derivation}ના કયા ભાગો બીજા કયા ભાગો પર આધાર રાખે છે તે
જોવાનું સહેલું બને છે.

આમ, પ્રથમ-ક્રમના તર્કશાસ્ત્ર જેવા કોઈ તર્કશાસ્ત્ર માટે જુદાં જુદાં
!!{derivation} તંત્રો !!a{sentence} \emph{પ્રમેય} હોવાનો અને કોઈ
બીજાં વાક્યોમાંથી !!a{sentence} !!{derivable} હોવાનો અર્થ જુદી
રીતે સ્પષ્ટ કરશે. આ કાર્ય સ્વયંસિદ્ધિમૂલક !!{derivation}s, સ્વાભાવિક
નિગમનો, સિક્વન્ટ !!{derivation}s, સત્ય-વૃક્ષો કે રિઝોલ્યૂશન વડે ખંડન
દ્વારા કરવામાં આવે, આપણે આ સંબંધો પ્રામાણ્ય અને ફલિતતાના અર્થાનુસારી
ખ્યાલો સાથે મળતા હોય તેમ ઇચ્છીએ છીએ. ``$!A$~પ્રમેય છે'' માટે
$\Proves !A$ અને ``$!A$ એ~$\Gamma$માંથી !!{derivable} છે'' માટે
``$\Gamma \Proves !A$'' લખીએ. $\Proves$ ગમે તે રીતે વ્યાખ્યાયિત હોય,
તે $\Entails$ સાથે નીચે પ્રમાણે મળવું જોઈએ:
\begin{enumerate}
\item $\Proves !A$ જો અને માત્ર જો $\Entails !A$
\item $\Gamma \Proves !A$ જો અને માત્ર જો $\Gamma \Entails !A$
\end{enumerate}
ઉપરની ``માત્ર જો'' દિશાને \emph{યથાર્થતા} કહે છે. !!^a{derivation}
તંત્ર યથાર્થ હોય તો !!{derivability} ફલિતતા (અથવા પ્રામાણ્ય)ની ખાતરી
આપે છે. દરેક યોગ્ય !!{derivation} તંત્ર યથાર્થ હોવું જ જોઈએ;
અયથાર્થ !!{derivation} તંત્રો કશા કામનાં નથી. આખરે,
!!a{derivation}નો સમગ્ર હેતુ
પ્રામાણ્ય અથવા ફલિતતાની વાક્યરચનાત્મક ખાતરી આપવાનો છે. આપણે રજૂ
કરીએ તે !!{derivation} તંત્રોની યથાર્થતા સાબિત કરીશું.

ઊલટી ``જો'' દિશા પણ મહત્વની છે: તેને \emph{પૂર્ણતા} કહે છે. પૂર્ણ
!!{derivation} તંત્ર એટલું સમર્થ હોય છે કે જ્યારે પણ $!A$ પ્રામાણ્ય
હોય ત્યારે તે $!A$ પ્રમેય છે તેમ અને જ્યારે પણ $\Gamma \Entails !A$
ત્યારે $\Gamma \Proves !A$ તેમ બતાવી શકે. પૂર્ણતા સ્થાપવી વધુ અઘરી
છે, અને કેટલાંક તર્કશાસ્ત્રો માટે કોઈ પૂર્ણ !!{derivation} તંત્ર હોતું
નથી. પ્રથમ-ક્રમના તર્કશાસ્ત્ર માટે આવું તંત્ર છે. કુર્ત ગેડલે 1929ના
પોતાના મહાનિબંધમાં પ્રથમ-ક્રમના તર્કશાસ્ત્ર માટેના !!a{derivation}
તંત્રની પૂર્ણતા સૌપ્રથમ સાબિત કરી હતી.

!!{derivation} તંત્રો સાથે જોડાયેલો બીજો ખ્યાલ \emph{સુસંગતતા} છે.
!!{sentence}sના કોઈ ગણમાંથી ગમે તે વસ્તુ !!{derive}d કરી શકાય તો તે
ગણને અસુસંગત અને નહિ તો સુસંગત કહે છે. અસુસંગતતા એ અસંતોષ્યતાનું
વાક્યરચનાત્મક પ્રતિરૂપ છે: અસંતોષ્ય ગણોની જેમ !!{sentence}sના
અસુસંગત ગણો સારા સિદ્ધાંતો બનતા નથી; તેમાં પાયાની ખામી હોય છે.
!!{sentence}sના સુસંગત ગણો સાચા કે ઉપયોગી હોય જ એવું નથી, પરંતુ
તેઓ તાર્કિક ઉપયોગિતાની આ લઘુતમ કસોટી તો પાર કરે છે. જુદાં
!!{derivation} તંત્રોમાં !!{sentence}sના ગણની સુસંગતતાની ચોક્કસ વ્યાખ્યા
જુદી હોઈ શકે, પરંતુ $\Proves$ની જેમ આપણે સુસંગતતા તેના અર્થાનુસારી
પ્રતિરૂપ—સંતોષ્યતા—સાથે મળતી હોય તેમ ઇચ્છીએ છીએ. $\Gamma$ સુસંગત હોય
જો અને માત્ર જો તે સંતોષ્ય હોય, એવું હંમેશાં હોવું જોઈએ. અહીં ``માત્ર
જો'' દિશા પૂર્ણતા જેટલી છે (સુસંગતતા સંતોષ્યતાની ખાતરી આપે છે), અને
``જો'' દિશા યથાર્થતા જેટલી છે (સંતોષ્યતા સુસંગતતાની ખાતરી આપે છે).
વાસ્તવમાં, શાસ્ત્રીય પ્રથમ-ક્રમના તર્કશાસ્ત્ર માટે યથાર્થતાનાં બંને
રૂપો પરસ્પર સમતુલ્ય છે, અને પૂર્ણતાનાં બંને રૂપો પણ પરસ્પર સમતુલ્ય છે.

\end{document}
